% Additional technical appendix. Root owns integration into main.tex.
\section{Continuous evaluation from partially observed episodes}
\label{app:asynchronous}

Waiting for the first unresolved episode can prevent a completed-prefix
analysis from using many later outcomes. Selecting only completed episodes
can instead introduce outcome-dependent selection. Here we retain the
enrollment order but replace each unresolved score by guaranteed bounds.
This is a worst-case completion construction: it does not require a model
for how quickly favorable and unfavorable outcomes become observable.
It specializes established full-data anytime inference and is closely
related to the pending-outcome stopping construction of
\citet{henzi2022forecast,henzi2022correction}.

\subsection{Latent enrollment order and the actual observation process}

Let $X_i$ be the complete bounded score of enrolled pair $i$, defined
under the frozen evaluation horizon and protocol. It may be the observed
oriented hierarchy score, its inverse-probability-weighted version from
Theorem~\ref{thm:pair_id}, or a component difference. The enrollment
filtration $(\mathcal F_i)$ used for the full-score argument need not be
the calendar-time filtration actually observed by the evaluator. We
require, rather than infer from the observation schedule, that
\[
 X_i\text{ is }\mathcal F_i\text{-measurable},\qquad
 a_i\leq X_i\leq b_i,\qquad
 \mu_i=E(X_i\mid\mathcal F_{i-1}),
\]
where $a_i,b_i$ are known predictable bounds. Put
$V_n=\sum_{i\leq n}(b_i-a_i)^2/4$ and
$\overline\mu_n=n^{-1}\sum_{i\leq n}\mu_i$.

At calendar time $t$, let $N(t)$ be the finite number of fully enrolled
and randomized pairs. The evaluator's information $\mathcal G_t$
contains whatever traces and labels have arrived by then. For each
$i\leq N(t)$, it supplies measurable bounds
$\ell_i(t),u_i(t)$ satisfying
\begin{equation}
 a_i\leq\ell_i(t)\leq X_i\leq u_i(t)\leq b_i
 \quad\text{simultaneously for every enrolled }i\text{ and every }t.
 \label{eq:async_enclosures}
\end{equation}
Bounds collapse to the score when it is determined. Outcome-dependent
delay, out-of-order updates, and partial information are allowed.
The observation process does not need to make $X_i$ independent of its
reveal time. For continuous calendar time, assume the relevant suprema
and events are measurable; an event-driven implementation with countably
many recorded updates suffices.

Condition~\eqref{eq:async_enclosures} is a simultaneous guarantee about
the \emph{frozen complete score}. A point prediction, an ordinary
pointwise 95\% interval, or an unvalidated judgment of likely success
does not satisfy it. If the guarantee instead holds on an event
$\mathcal A$ with $P(\mathcal A)\geq1-\eta$, all the error bounds below
increase by at most $\eta$, without requiring independence.

\paragraph{A sufficient experimental model.}
One transparent model samples independent full potential episode records
and their partial-reveal paths for each pair, uses fixed candidate systems
and a fixed outcome protocol, and has no cross-pair interference. The
record for pair $i$ can include arbitrarily dependent final outcomes and
reveal times. Pair orientation is then drawn with a fresh recorded
randomization coin; its probability can depend on information from
previously enrolled pairs. If that information is a function of their
full records, it is measurable in the filtration revealing those records
in enrollment order. The pair-score identification and bounded-score
arguments apply in this filtration. Balanced orientation with a fixed
independent workload supplies the stationary case. In a design-based
version, full potential records and reveal paths can be conditioned upon,
with independent assignment coins supplying the randomness.

These are sufficient conditions, not a claim that any asynchronous
production system satisfies them. In particular, revealing all earlier
final records in a mathematical filtration can change a conditional-mean
restriction. Shared infrastructure effects, interference, or adaptive
candidate updates require their own justification. The hierarchy must
not be refitted after a pending episode's partial outcome is seen.
Time-uniform coverage and threshold-crossing guarantees transfer by
pathwise event inclusion; free confidence-sequence and p-process lifting
are established in \citet{choe2024filtrations}. Their delayed-forecast
discussion is also a close structural precedent. Such a transfer does
not establish a calendar-time e-process or the required latent
enrollment-order conditional-mean restriction.

\subsection{Confidence sequences from all enrolled pairs}

Let $B_\alpha(v)$ be the normal-mixture boundary in
Theorem~\ref{thm:normal_cs}. For any $1\leq n\leq N(t)$ define
\begin{equation}
 L_{n,t}=\frac{\sum_{i=1}^n\ell_i(t)-B_\alpha(V_n)}n,
 \qquad
 U_{n,t}=\frac{\sum_{i=1}^n u_i(t)+B_\alpha(V_n)}n.
 \label{eq:async_cs}
\end{equation}
In particular, $n=N(t)$ uses every enrolled pair, including unresolved
pairs before or after a completed pair. The bounds use $n$, not the
number of completed episodes, in both the sum and the boundary.

\begin{theorem}[Pathwise transfer of the full-score confidence sequence]
\label{thm:async_cs}
Under the full-score assumptions of Theorem~\ref{thm:normal_cs} and
the simultaneous enclosures \eqref{eq:async_enclosures},
\[
 P\{L_{n,t}\leq\overline\mu_n\leq U_{n,t}
       \text{ for every }t\text{ and every }1\leq n\leq N(t)\}
 \geq1-\alpha.
\]
If the enclosures hold only on $\mathcal A$, the lower bound is
$1-\alpha-\eta$.
\end{theorem}
\begin{proof}
Let $\mathcal C$ be the event of full-score coverage for every integer
$n$, which has probability at least $1-\alpha$ by
Theorem~\ref{thm:normal_cs}. On $\mathcal C$, for every $n$,
\[
 \frac{\sum_{i\leq n}X_i-B_\alpha(V_n)}n
 \leq\overline\mu_n\leq
 \frac{\sum_{i\leq n}X_i+B_\alpha(V_n)}n.
\]
On the enclosure event, replacing the left sum by its lower bound and
the right sum by its upper bound only widens this interval. This is a
pathwise statement for every available $(n,t)$, so there is no need to
assert that $N(t)$ or a display time is an $\mathcal F$-stopping time.
The desired event contains $\mathcal C$ when the enclosures are certain,
and contains $\mathcal C\cap\mathcal A$ otherwise. A union bound gives
the second probability statement.
\end{proof}

The width has an exact decomposition:
\begin{equation}
 U_{n,t}-L_{n,t}
 =\frac{2B_\alpha(V_n)}n+
   \frac1n\sum_{i=1}^n\{u_i(t)-\ell_i(t)\}.
 \label{eq:async_width}
\end{equation}
The second term is the price of unresolved outcome information. If
$X_i\in[-B,B]$ and only $M$ of the first $n$ scores remain unresolved,
its contribution is at most $2BM/n$ when the enclosures are clipped to
$[-B,B]$. Thus one unresolved early pair
need not prevent learning from arbitrarily many later pairs.
This is a bound on information loss, not a guarantee that the method
always produces a narrower interval than waiting for a smaller completed
prefix; those procedures can also concern different running targets.

\subsection{Partial-score betting evidence}

For every fixed-$B$ result below, clip the predictable bounds and all
enclosures to $[-B,B]$, so that $-B\leq a_i\leq\ell_i(t)\leq
X_i\leq u_i(t)\leq b_i\leq B$. Clipping preserves containment.
For clarity use a fixed bound $X_i\in[-B,B]$ and a fixed gate threshold
$c\in(-B,B)$. Choose $J$ stakes
$0\leq\lambda_k<1/(B+c)$ and weights $w_k\geq0$ with
$\sum_kw_k=1$ before viewing the stream. The full-score mixture is
\[
 E_n^{\mathrm{full}}(c)=
 \sum_{k=1}^{J}w_k\prod_{i=1}^n\{1+\lambda_k(X_i-c)\}.
\]
The observable partial-score lower wealth for any enrolled prefix is
\begin{equation}
 \underline E_{n,t}(c)=
 \sum_{k=1}^{J}w_k\prod_{i=1}^n
                      \{1+\lambda_k(\ell_i(t)-c)\}.
 \label{eq:async_wealth}
\end{equation}
Past factors are recomputed or updated when an enclosure improves; they
are not counted as fresh independent observations. Stakes and mixture
weights remain the ones committed before seeing this stream.

\begin{theorem}[Threshold validity of partial-score lower wealth]
\label{thm:async_betting}
If $E(X_i\mid\mathcal F_{i-1})\leq c$ for every $i$, then
\[
 P\left\{\exists t,\ 1\leq n\leq N(t):
             \underline E_{n,t}(c)\geq1/\alpha\right\}\leq\alpha.
\]
If the enclosures hold only on $\mathcal A$, the bound is $\alpha+\eta$.
This establishes a valid anytime threshold test. It does not assert that
$\underline E_{N(t),t}(c)$ is a martingale or an e-process in
$(\mathcal G_t)$.
\end{theorem}
\begin{proof}
Under the null, each full-score factor is positive and has conditional
expectation at most one. Hence $E_n^{\mathrm{full}}(c)$ is a
nonnegative supermartingale starting at one, by the same argument as
Theorem~\ref{thm:betting}. On the enclosure event, nonnegative stakes
give
\[
 0<1+\lambda_k(\ell_i(t)-c)
 \leq1+\lambda_k(X_i-c).
\]
Products and their nonnegative mixtures preserve the inequality, so
$\underline E_{n,t}(c)\leq E_n^{\mathrm{full}}(c)$ for all $(n,t)$.
Every partial-wealth threshold crossing therefore implies a crossing of
the full-score process at some integer $n$. Ville's inequality bounds
the latter event by $\alpha$. A possible enclosure failure adds at most
$\eta$. Nothing in this event inclusion proves a calendar-time
conditional-expectation identity, which explains the final qualification.
\end{proof}

If desired, the process
\[
 p_t=\min\!\left\{1,
 \left[\sup_{s\leq t}\ \max_{1\leq n\leq N(s)}
                         \underline E_{n,s}(c)\right]^{-1}\right\}
\]
is an anytime-valid p-value process under exact enclosures, with an empty
maximum interpreted as one. The threshold statement follows directly
from Theorem~\ref{thm:async_betting}. Calling its reciprocal an e-process
would require an additional expectation guarantee and is not warranted
by this argument alone.

\begin{corollary}[A prefix envelope contains the completed-prefix test]
\label{cor:async_envelope}
\leavevmode\par
Define $\underline E_t^{\max}(c)
=\max_{1\leq n\leq N(t)}\underline E_{n,t}(c)$. It has the same
threshold guarantee as Theorem~\ref{thm:async_betting}, without dividing
$\alpha$ over prefixes. If $m(t)$ is a completely resolved enrollment
prefix with $\ell_i(t)=u_i(t)=X_i$ for $i\leq m(t)$, then
\[
 \underline E_t^{\max}(c)\geq E_{m(t)}^{\mathrm{full}}(c).
\]
Consequently the envelope test cannot cross later than the
completed-prefix test with the same full process and monitoring times.
\end{corollary}
\begin{proof}
The threshold event is exactly a subset of the event already bounded in
Theorem~\ref{thm:async_betting}. At the resolved prefix $m(t)$, partial
and full factors coincide, and this prefix is one of the maximands.
The crossing-time ordering follows pointwise.
\end{proof}

This corollary concerns the envelope, not merely the all-enrolled
quantity $\underline E_{N(t),t}$. Adding many unresolved pairs can reduce
that latter quantity, so the all-enrolled implementation has no universal
speed ordering relative to completed-prefix inference. If each enclosure
narrows over calendar time, a fixed-prefix lower wealth is nondecreasing
as information arrives. Revisions that widen a previously overconfident
enclosure indicate that its claimed pathwise guarantee was not valid.

For stationary means, the maximum over enrolled prefixes of $L_{n,t}$
in \eqref{eq:async_cs} is also a valid lower bound on the same constant
$\mu$. Under drift, $\overline\mu_n$ depends on $n$, so maximizing
lower bounds across prefixes does not produce a lower bound for
$\overline\mu_{N(t)}$. The drifting analysis should retain the selected
prefix and its corresponding target explicitly.

\subsection{Constructing guaranteed hierarchy enclosures from traces}

Let $\mathcal C_i^A(t)$ and $\mathcal C_i^B(t)$ contain all final
episode records consistent with the information observed for the assigned
$A$ and $B$ episodes. Require that the actual complete records belong
to these sets simultaneously. For the unweighted hierarchy define
\begin{equation}
 \ell_i^h(t)=\inf_{a\in\mathcal C_i^A(t),\ b\in\mathcal C_i^B(t)}h(a,b),
 \qquad
 u_i^h(t)=\sup_{a\in\mathcal C_i^A(t),\ b\in\mathcal C_i^B(t)}h(a,b).
 \label{eq:async_kernel_set}
\end{equation}
These are bounds on the observed oriented comparison, not on an
unobserved alternative-arm outcome for the same user. Measurability of
the bounds is required; finite outcome-state enumeration is sufficient
for binary tiers and possible signs of the resource comparison. The
true final records are among the feasible completions, proving
\eqref{eq:async_enclosures}. A superset of feasible completions gives
conservative bounds even if the separate completion sets ignore some
within-pair constraints.

For a component score $g(a)-g(b)$, if the traces guarantee
$g(a)\in[g_A^-,g_A^+]$ and $g(b)\in[g_B^-,g_B^+]$, use
\[
 \ell_i^g=g_A^- -g_B^+,\qquad
 u_i^g=g_A^+ -g_B^-.
\]
For the inverse-probability pair score, multiply both endpoints by the
known positive factor $1/(2q_i)$ or $1/\{2(1-q_i)\}$ corresponding to
the realized orientation. The resulting global or predictable score
range must also be used in the full-score bound.

The hierarchy gives useful early conclusions without requiring every
endpoint to be finalized. If all higher-priority tiers are guaranteed
tied and the next tier is guaranteed decisive, its sign is final even
when lower-priority resource records remain open. A stronger implication
can sometimes combine several possible completions. For example, use
success then cost with joint failures tied. Suppose $A$ has completed
with verified success at cost 1, while $B$ is unresolved but has already
incurred cost 2. Under a 5\% relative cost tolerance and nondecreasing
recorded cost, $A$ wins if $B$ ultimately fails and also wins if $B$
ultimately succeeds: any final $B$ cost is at least 2. Therefore
$\ell_i^h=u_i^h=1$ before $B$ completes. The success-difference enclosure
is only $[0,1]$, so superiority in success has not been established by
that episode. A higher-priority safety tier still unresolved would
invalidate this particular early win certificate.

Absence of a violation so far is not a certificate that the final
episode is violation-free. Similarly, accumulated cost is a lower bound
only when the recorded cost cannot later decrease. Refunds, post hoc
verifier changes, or continuing activity after an apparently successful
termination require broader completion sets. If a partial trace is
missing, its possible completions remain in the bounds. Evaluation
cannot improve evidence by deleting that pair.

\subsection{When can partial information eventually suffice?}

\begin{proposition}[Vanishing unresolved information]
\label{prop:async_consistency}
Suppose $X_i\in[-B,B]$ have fixed conditional mean $\mu$, and use a
fixed normal-mixture tuning parameter. Along any observed sequence of
calendar looks $t_k$ with $n_k=N(t_k)\to\infty$, suppose
\[
 \frac1{n_k}\sum_{i\leq n_k}
       \{u_i(t_k)-\ell_i(t_k)\}\longrightarrow0
 \quad\text{almost surely}.
\]
Then both bounds $L_{n_k,t_k}$ and $U_{n_k,t_k}$ converge to $\mu$
almost surely. Thus a finite collection of stationary criteria with
strict positive gaps above their thresholds eventually passes the
partial-score confidence-bound rule, provided this condition holds
for every criterion.
\end{proposition}
\begin{proof}
Bounded martingale differences give
$n^{-1}\sum_{i\leq n}(X_i-\mu)\to0$ almost surely. One direct proof
uses the exponential fixed-time Hoeffding bound, its summability for
any fixed positive deviation, and Borel--Cantelli, followed by a
countable sequence of deviations tending to zero. By enclosure,
\[
 0\leq\frac1{n_k}\sum_{i\leq n_k}\{X_i-\ell_i(t_k)\}
 \leq\frac1{n_k}\sum_{i\leq n_k}\{u_i(t_k)-\ell_i(t_k)\},
\]
and likewise for $u_i(t_k)-X_i$. Hence the averages of lower and upper
endpoints converge to $\mu$. The normal-mixture radius tends to zero
because $V_n\leq nB^2$. The finite-gate conclusion follows by taking
the maximum of the almost surely finite eventual-passing indices.
\end{proof}

The condition permits an unresolved first pair forever if the total
unresolved width is $o(n)$. It can also hold when complete traces remain
open but their hierarchy signs have already been certified. Conversely,
if all scores remain enclosed only by $[-B,B]$, the lower bound cannot
establish a positive mean without additional information. Robustness to
informative delay is obtained through this explicit uncertainty cost.

\begin{proposition}[A bound on the loss of log wealth]
\label{prop:async_log_penalty}
For a fixed positive stake $\lambda<1/(B+c)$ let
$\gamma=1-\lambda(B+c)>0$. For any $n\leq N(t)$, the corresponding
full and partial single-stake log products satisfy
\[
 0\leq\log E^{\mathrm{full},\lambda}_n
          -\log\underline E^\lambda_{n,t}
 \leq\frac\lambda\gamma
              \sum_{i=1}^n\{X_i-\ell_i(t)\}.
\]
If the full-score sequence is i.i.d., the stake has strictly positive
expected log growth, and the average upper-minus-lower width vanishes
along $n_k=N(t_k)\to\infty$, its partial log product divided by $n_k$
has the same positive limit. A fixed mixture assigning positive weight
to that stake therefore eventually crosses every finite threshold.
\end{proposition}
\begin{proof}
The function $f(x)=\log\{1+\lambda(x-c)\}$ is increasing on $[-B,B]$
and has derivative at most $\lambda/\gamma$. Apply the mean value
theorem to each $f(X_i)-f(\ell_i(t))$ and sum. The strong law gives the
full log-growth limit; the normalized difference vanishes by the width
assumption and the displayed inequality. A fixed positive mixture
weight contributes only a constant to the logarithm of its term.
\end{proof}

\subsection{Deployment conjunctions and the scope of the guarantee}

Under fixed conditional means, every required stationary gate can use
its partial-score threshold test at level $\alpha$. If one fixed gate
null is true, an eventual deployment requires that gate's partial wealth
to cross, whose probability is at most $\alpha$ by
Theorem~\ref{thm:async_betting}. Gates may cross at different times or
prefix lengths because they concern fixed stationary parameters. This
is the same intersection--union argument as Theorem~\ref{thm:iut},
and does not require independence between reveal times or criteria.
An enclosure failure event common to all gates adds at most $\eta$.

For arbitrary drifting running means, use the bounds
\eqref{eq:async_cs} for the same stated enrolled prefix, allocate the
confidence-sequence error across all required criteria, and require
every lower bound to exceed its threshold. The union-bound proof of
Theorem~\ref{thm:drift_gate} then applies pathwise. A historical
partial-bound crossing cannot be retained as evidence about a different
current drifting target. The same candidate-selection and repeated-test
restrictions continue to apply.

These results address eventual outcomes through a \emph{fixed episode
horizon}. They are distinct from the calendar-time cumulative reward
target studied by \citet{lindon2026delayed}, who analyze delayed
randomized outcomes and their filtration difficulties. No assertion of
calendar-time martingale structure is needed here; nor do our bounds
contradict that work's obstruction for a different, calendar-time
treatment-effect process. Stopping enrollment does not authorize
changing the endpoint definition for pending episodes. If deployment
terminates pending runs, either the prespecified intervention includes
that termination or the estimand remains the stated hypothetical
continuation through the frozen horizon.
